\ifdefined\gkver\else\def\gkver{student}\fi
\documentclass[\gkver]{gaokaozhenti}
\begin{document}

\chapter{2014年全国理综大纲}
%% paperType: 理综物理部分

\begin{enumerate}
\item
%% number: 14
%% paperName: 2014年全国理综大纲第 14 题 6 分
%% typeId: 1
%% type: 单选题
%% chapter: 直线运动
%% point: 运动图象
%% method: 无
%% score: 6
%% degree: 650
%% duplicateId: 0
%% body:
一质点沿 $x$ 轴做直线运动，其 $v-t$ 图像如图所示。质点在 $t=0$ 时位于 $x=5\Um$ 处，开始沿 $x$ 轴正向运动。当 $t=8\Us$ 时，质点在 $x$ 轴上的位置为\xzanswer{B}
\onepicture[w=6.0cm]{figs/14.png}
\fourchoices[answer=B]
{$x=3\Um$}
{$x=8\Um$}
{$x=9\Um$}
{$x=14\Um$}
%% answer: B
%% memo:
\memoanswer{
根据图像表示的物理意义可知，图线与时间轴围成的面积表示物体的位移，面积在时间轴之上，表示位移为正，反之表示位移为负。由图像可知 $8\Us$ 内质点的位移为 $s=\frac{(2+4)\times2}{2}\Um-\frac{(2+4)\times1}{2}\Um=3\Um$，又因为初始时刻质点的位置为 $x=5\Um$ 处，所以 $8\Us$ 末质点在 $8\Um$ 处，选项 B 正确。
}

\item
%% number: 15
%% paperName: 2014年全国理综大纲第 15 题 6 分
%% typeId: 1
%% type: 单选题
%% chapter: 电场
%% point: 带电粒子的直线运动
%% method: 无
%% score: 6
%% degree: 650
%% duplicateId: 0
%% body:
地球表面附近某区域存在大小为 $150\UNC$、方向竖直向下的电场。一质量为 $1.00\times10^{-4}\Ukg$、带电量为 $-1.00\times10^{-7}\UC$ 的小球从静止释放，在电场区域内下落 $10.0\Um$。对此过程，该小球的电势能和动能的改变量分别为（重力加速度大小取 $9.80\Umsq$，忽略空气阻力）\xzanswer{D}
\fourchoices[answer=D]
{$-1.50\times10^{-4}\UJ$ 和 $9.95\times10^{-3}\UJ$}
{$1.50\times10^{-4}\UJ$ 和 $9.95\times10^{-3}\UJ$}
{$-1.50\times10^{-4}\UJ$ 和 $9.65\times10^{-3}\UJ$}
{$1.50\times10^{-4}\UJ$ 和 $9.65\times10^{-3}\UJ$}
%% answer: D
%% memo:
\memoanswer{
设小球下落的高度为 $h$，则电场力做的功 $W_1=-qEh=-1.5\times10^{-4}\UJ$，电场力做负功，电势能增加，所以电势能增加 $1.5\times10^{-4}\UJ$；重力做的功 $W_2=mgh=9.8\times10^{-3}\UJ$，合力做的功 $W=W_1+W_2=9.65\times10^{-3}\UJ$，根据动能定理可知 $\Delta E_k=W=9.65\times10^{-3}\UJ$，因此选项 D 正确。
}

\item
%% number: 16
%% paperName: 2014年全国理综大纲第 16 题 6 分
%% typeId: 2
%% type: 多选题
%% chapter: 气体性质
%% point: 理想气体状态方程
%% method: 无
%% score: 6
%% degree: 450
%% duplicateId: 0
%% body:
对于一定量的稀薄气体，下列说法正确的是\xzanswer{BD}
\fourchoices[answer=BD]
{压强变大时，分子热运动必然变得剧烈}
{保持压强不变时，分子热运动可能变得剧烈}
{压强变大时，分子间的平均距离必然变小}
{压强变小时，分子间的平均距离可能变小}
%% answer: BD
%% memo:
\memoanswer{
一定质量的稀薄气体可以看做理想气体，分子运动的剧烈程度与温度有关，温度越高，分子运动的越剧烈；压强变大可能是的原因是体积变小或温度升高，所以压强变大，分子热运动不一定剧烈，AC 项错误，B 项正确；压强变小时，也可能体积不变，可能变大，也可能变小；温度可能降低，可能不变，可能升高，所以分子间距离不能确定，D 项正确。
}

\item
%% number: 17
%% paperName: 2014年全国理综大纲第 17 题 6 分
%% typeId: 1
%% type: 单选题
%% chapter: 光学
%% point: 光的干涉
%% method: 无
%% score: 6
%% degree: 450
%% duplicateId: 0
%% body:
在双缝干涉实验中，一钠灯发出的波长为 $589\Unm$ 的光，在距双缝 $1.00\Um$ 的屏上形成干涉图样。图样上相邻两明纹中心间距为 $0.350\Ucm$，则双缝的间距为\xzanswer{C}
\fourchoices[answer=C]
{$2.06\times10^{-7}\Um$}
{$2.06\times10^{-4}\Um$}
{$1.68\times10^{-4}\Um$}
{$1.68\times10^{-3}\Um$}
%% answer: C
%% memo:
\memoanswer{
双缝干涉相邻条纹间距 $\Delta x=\frac{l}{d}\lambda$，则 $d=\frac{l}{\Delta x}\lambda=1.68\times10^{-4}\Um$，C 正确。
}

\item
%% number: 18
%% paperName: 2014年全国理综大纲第 18 题 6 分
%% typeId: 1
%% type: 单选题
%% chapter: 机械波
%% point: 干涉衍射
%% method: 无
%% score: 6
%% degree: 450
%% duplicateId: 0
%% body:
两列振动方向相同、振幅分别为 $A_1$ 和 $A_2$ 的相干简谐横波相遇。下列说法正确的是\xzanswer{AD}
\fourchoices[answer=AD]
{波峰与波谷相遇处质点的振幅为 $|A_1-A_2|$}
{波峰与波峰相遇处质点离开平衡位置的位移始终为 $A_1+A_2$}
{波峰与波谷相遇处质点的位移总是小于波峰与波峰相遇处质点的位移}
{波峰与波峰相遇处质点的振幅一定大于波峰与波谷相遇处质点的振幅}
%% answer: AD
%% memo:
\memoanswer{
根据波的叠加原理，两列波相遇互不干扰，所以波峰与波谷相遇，质点的振幅变为 $|A_1-A_2|$，A 项正确；任何一点质点都不会总处于波峰或波谷，总是在平衡位置附近往复运动，BC 项错误；波峰与波峰相遇处的质点处于振动加强点，所以振幅较大，波谷与波峰处的质点处于振动减弱点，振幅偏小，D 项正确。
}

\item
%% number: 19
%% paperName: 2014年全国理综大纲第 19 题 6 分
%% typeId: 1
%% type: 单选题
%% chapter: 运动和力
%% point: 牛顿定律的应用
%% method: 无
%% score: 6
%% degree: 450
%% duplicateId: 0
%% body:
一物块沿倾角为 $\theta$ 的斜坡向上滑动。当物块的初速度为 $v$ 时，上升的最大高度为 $H$，如图所示；当物块的初速度为 $\frac{v}{2}$ 时，上升的最大高度记为 $h$。重力加速度大小为 $g$。物块与斜坡间的动摩擦因数和 $h$ 分别为\xzanswer{D}
\onepicture[w=5.5cm]{figs/19.png}
\fourchoices[answer=D]
{$\tan\theta$ 和 $\frac{H}{2}$}
{$\left(\frac{v^2}{2gH}-1\right)\tan\theta$ 和 $\frac{H}{2}$}
{$\tan\theta$ 和 $\frac{H}{4}$}
{$\left(\frac{v^2}{2gH}-1\right)\tan\theta$ 和 $\frac{H}{4}$}
%% answer: D
%% memo:
\memoanswer{
物块向上做匀减速直线运动，根据受力可知，其加速度为 $a=g\sin\theta+\mu g\cos\theta$；再根据匀变速位移速度关系，有 $2(g\sin\theta+\mu g\cos\theta)H=v^2$，$2(g\sin\theta+\mu g\cos\theta)h=\left(\frac{v}{2}\right)^2$，联立解得 $\mu=\left(\frac{v^2}{2gH}-1\right)\tan\theta$，$h=\frac{H}{4}$，D 项正确。
}

\item
%% number: 20
%% paperName: 2014年全国理综大纲第 20 题 6 分
%% typeId: 1
%% type: 单选题
%% chapter: 电磁感应
%% point: 楞次定律推广
%% method: 无
%% score: 6
%% degree: 650
%% duplicateId: 0
%% body:
很多相同的绝缘铜圆环沿竖直方向叠放，形成一很长的竖直圆筒。一条形磁铁沿圆筒的中心轴竖直放置，其下端与圆筒上端开口平齐。让条形磁铁从静止开始下落。条形磁铁在圆筒中的运动速率\xzanswer{C}
\fourchoices[answer=C]
{均匀增大}
{先增大，后减小}
{逐渐增大，趋于不变}
{先增大，再减小，最后不变}
%% answer: C
%% memo:
\memoanswer{
小磁铁在圆环中下落的过程由于穿过圆环的磁通量发生变化，所以在圆环中会感应电流阻碍磁铁的下落，又因为圆筒很长，所以当磁铁受到的阻力等于重力时，其速度保持趋于不变，选项 C 正确。
}

\item
%% number: 21
%% paperName: 2014年全国理综大纲第 21 题 6 分
%% typeId: 1
%% type: 单选题
%% chapter: 运动和力
%% point: 动量守恒定律
%% method: 无
%% score: 6
%% degree: 450
%% duplicateId: 0
%% body:
一中子与一质量数为 $A$（$A>1$）的原子核发生弹性正碰。若碰前原子核静止，则碰撞前与碰撞后中子的速率之比为\xzanswer{A}
\fourchoices[answer=A]
{$\frac{A+1}{A-1}$}
{$\frac{A-1}{A+1}$}
{$\frac{4A}{(A+1)^2}$}
{$\frac{(A+1)^2}{(A-1)^2}$}
%% answer: A
%% memo:
\memoanswer{
设碰撞前后中子的速率分别为 $v_1$、$v'_1$，碰撞后原子核的速率为 $v_2$，中子的质量为 $m_1$，原子核的质量为 $m_2$，则 $m_2=Am_1$。根据完全弹性碰撞规律可得 $m_1v_1=m_2v_2+m_1v'_1$，$\frac{1}{2}m_1v_1^2=\frac{1}{2}m_2v_2^2+\frac{1}{2}m_1v'^2_1$，解得碰后中子的速率 $v'_1=\left|\frac{m_1-m_2}{m_1+m_2}\right|v_1=\frac{A-1}{A+1}v_1$，因此碰撞前后中子速率之比 $\frac{v_1}{v'_1}=\frac{A+1}{A-1}$，选项 A 正确。
}

\item
%% number: 22
%% paperName: 2014年全国理综大纲第 22 题 6 分
%% typeId: 4
%% type: 实验题
%% chapter: 直线运动
%% point: 打点计时器公式
%% method: 无
%% score: 6
%% degree: 650
%% duplicateId: 0
%% body:
现用频闪照相方法来研究物块的变速运动。在一小物块沿斜面向下运动的过程中，用频闪相机拍摄的不同时刻物块的位置如图所示。拍摄时频闪频率是 $10\UHz$；通过斜面上固定的刻度尺读取的 $5$ 个连续影像间的距离依次为 $x_1$、$x_2$、$x_3$、$x_4$。已知斜面顶端的高度 $h$ 和斜面的长度 $s$。数据如下表所示（单位：$\Ucm$）。重力加速度大小 $g=9.80\Umsq$。

\begin{center}
\begin{tabular}{|c|c|c|c|c|c|}
\hline
$x_1$ & $x_2$ & $x_3$ & $x_4$ & $h$ & $s$ \\ \hline
10.76 & 15.05 & 19.34 & 23.65 & 48.00 & 80.00 \\ \hline
\end{tabular}
\end{center}

根据表中数据，完成下列填空：
\begin{enumerate}
	\item
	物块的加速度 $a=$\tkanswer[1.8cm]{4.30}$\Umsq$（保留 3 位有效数字）。
	\item
	因为\tkanswer[6.5cm]{物块加速度小于 $g\frac{h}{s}=5.88\Umsq$}，可知斜面是粗糙的。
\end{enumerate}
\onepicture[align=right, w=8.0cm]{figs/22.png}
%% answer:
\jdanswer{
\begin{enumerate}
	\item
	4.30（填“4.29”或“4.31”同样给分）

	\item
	物块加速度小于 $g\frac{h}{s}=5.88\Umsq$（或：物块加速度小于物块沿光滑斜面下滑的加速度）

\end{enumerate}
}
%% memo:
\memoanswer{
（1）根据逐差法求出加速度 $a=\frac{(x_3+x_4)-(x_1+x_2)}{(2T)^2}=4.30\Umsq$。

（2）根据牛顿第二定律，物块沿光滑斜面下滑的加速度 $a'=g\sin\theta=g\frac{h}{s}=5.88\Umsq$，由于 $a<a'$，可知斜面是粗糙的。
}

\item
%% number: 23
%% paperName: 2014年全国理综大纲第 23 题 12 分
%% typeId: 4
%% type: 实验题
%% chapter: 电路
%% point: 测电动势与内阻
%% method: 无
%% score: 12
%% degree: 450
%% duplicateId: 0
%% body:
现要测量某电源的电动势和内阻。可利用的器材有：电流表 A，内阻为 $1.00\UO$；电压表 V；阻值未知的定值电阻 $R_1$、$R_2$、$R_3$、$R_4$、$R_5$；开关 $S$；一端连有鳄鱼夹 $P$ 的导线 1，其他导线若干。某同学设计的测量电路如图（a）所示。

\onepicture[showlabel=false, w=6.0cm]{figs/23a.png}

\onepicture[showlabel=false, w=9.0cm]{figs/23b.png}

\begin{enumerate}
	\item
	按图（a）在实物图（b）中画出连线，并标出导线 1 和其 $P$ 端。
	\item
	测量时，改变鳄鱼夹 $P$ 所夹的位置，使 $R_1$、$R_2$、$R_3$、$R_4$、$R_5$ 依次串入电路，记录对应的电压表的示数 $U$ 和电流表的示数 $I$。数据如下表所示。根据表中数据，在图（c）中的坐标纸上将所缺数据点补充完整，并画出 $U-I$ 图线。

	\begin{center}
	\begin{tabular}{|c|c|c|c|c|c|}
	\hline
	$I$（$\UmA$） & 193 & 153 & 111 & 69 & 30 \\ \hline
	$U$（$\UV$） & 2.51 & 2.59 & 2.68 & 2.76 & 2.84 \\ \hline
	\end{tabular}
	\end{center}

	\onepicture[showlabel=false, w=7.5cm]{figs/23c.png}

	\item
	根据 $U-I$ 图线求出电源的电动势 $E=$\tkanswer[1.6cm]{2.90}$\UV$，内阻 $r=$\tkanswer[1.6cm]{1.03}$\UO$。
\end{enumerate}
%% answer:
\jdanswer{
\begin{enumerate}
	\item
	连线如图所示。

	\item
	$U-I$ 图线如图所示。

	\item
	$E=2.90\UV$，$r=1.03\UO$

\end{enumerate}
}
%% memo:
\memoanswer{
（1）电流表的内阻已知，所以对实验数据的测量没有误差，可以将其看做电源的内阻；用固定电阻代替滑动变阻器，可以知道外电路的电阻，根据闭合电路欧姆定律可以测量出电源的电动势和内阻，实物连接如图：

\onepicture[showlabel=false, w=7.0cm]{figs/23_an1.png}

（2）根据给出的数据进行描点，用直线连接即可：

\onepicture[showlabel=false, w=7.0cm]{figs/23_an2.png}

（3）根据 $U=E-I(r_A+r)$，图线的斜率表示电源的内阻和电流表的内阻之和，由图像可知斜率为：$r=k-r_A=\frac{\Delta U}{\Delta I}-r_A=1.03\UO$，纵截距表示电源的电动势，即 $E=2.90\UV$。
}

\item
%% number: 24
%% paperName: 2014年全国理综大纲第 24 题 12 分
%% typeId: 6
%% type: 计算题
%% chapter: 运动和力
%% point: 动量守恒定律
%% method: 无
%% score: 12
%% degree: 650
%% duplicateId: 0
%% body:
冰球运动员甲的质量为 $80.0\Ukg$。当他以 $5.0\Ums$ 的速度向前运动时，与另一质量为 $100\Ukg$、速度为 $3.0\Ums$ 的迎面而来的运动员乙相撞。碰后甲恰好静止。假设碰撞时间极短，求：
\begin{enumerate}
	\item
	碰后乙的速度的大小；
	\item
	碰撞中总机械能的损失。
\end{enumerate}
%% answer:
\jdanswer{
\begin{enumerate}
	\item
	$1.0\Ums$

	\item
	$1400\UJ$

\end{enumerate}
}
%% memo:
\memoanswer{
（1）设运动员甲、乙的质量分别为 $m$、$M$，碰前速度大小分别为 $v$、$V$，碰后乙的速度大小为 $V'$。由动量守恒定律有

$mv-MV=MV'$ ①

代入数据得

$V'=1.0\Ums$ ②

（2）设碰撞过程中总机械能的损失为 $\Delta E$，应有

$\frac{1}{2}mv^2+\frac{1}{2}MV^2=\frac{1}{2}MV'^2+\Delta E$ ③

联立②③式，代入数据得

$\Delta E=1400\UJ$ ④
}

\item
%% number: 25
%% paperName: 2014年全国理综大纲第 25 题 20 分
%% typeId: 1
%% type: 单选题
%% chapter: 磁场
%% point: 带电粒子在复合场中的运动
%% method: 无
%% score: 20
%% degree: 450
%% duplicateId: 0
%% body:
如图，在第一象限存在匀强磁场，磁感应强度方向垂直于纸面（$xy$ 平面）向外；在第四象限存在匀强电场，方向沿 $x$ 轴负向。在 $y$ 轴正半轴上某点以与 $x$ 轴正向平行、大小为 $v_0$ 的速度发射出一带正电荷的粒子，该粒子在（$d$，$0$）点沿垂直于 $x$ 轴的方向进入电场。不计重力。若该粒子离开电场时速度方向与 $y$ 轴负方向的夹角为 $\theta$，求：
\begin{enumerate}
	\item
	电场强度大小与磁感应强度大小的比值；
	\item
	该粒子在电场中运动的时间。
\end{enumerate}
\onepicture[align=right, w=4.5cm]{figs/25.png}
%% answer:
\jdanswer{
\begin{enumerate}
	\item
	$\frac{E}{B}=\frac{1}{2}v_0\tan^2\theta$

	\item
	$t=\frac{2d}{v_0\tan\theta}$

\end{enumerate}
}
%% memo:
\memoanswer{
（1）如图粒子进入磁场后做匀速圆周运动，设磁感应强度大小为 $B$，粒子质量与所带电荷量分别为 $m$ 和 $q$，圆周运动的半径为 $R_0$，由洛伦兹力公式及牛顿第二定律得

$qv_0B=\frac{mv_0^2}{R_0}$

由题给条件和几何关系可知 $R_0=d$

设电场强度大小为 $E$，粒子进入电场后沿 $x$ 轴负方向的加速度大小为 $a_x$，在电场中运动的时间为 $t$，离开电场时沿 $x$ 轴负方向的速度大小为 $v_y$。由牛顿定律及运动学公式得

$qE=ma_x$

$v_x=at$

$\frac{v_x}{2}=d$

粒子在电场中做类平抛运动，如图所示

$\tan\theta=\frac{v_y}{v_0}$

\onepicture[showlabel=false, w=4.0cm]{figs/25_an.png}

联立得 $\frac{E}{B}=\frac{1}{2}v_0\tan^2\theta$

（2）同理可得 $t=\frac{2d}{v_0\tan\theta}$
}

\item
%% number: 26
%% paperName: 2014年全国理综大纲第 26 题 22 分
%% typeId: 6
%% type: 计算题
%% chapter: 曲线运动
%% point: 人造地球卫星
%% method: 无
%% score: 22
%% degree: 250
%% duplicateId: 0
%% body:
已知地球的自转周期和半径分别为 $T$ 和 $R$，地球同步卫星 A 的圆轨道半径为 $h$。卫星 B 沿半径为 $r$（$r<h$）的圆轨道在地球赤道的正上方运行，其运行方向与地球自转方向相同。求：
\begin{enumerate}
	\item
	卫星 B 做圆周运动的周期；
	\item
	卫星 A 和 B 连续地不能直接通讯的最长时间间隔（信号传输时间可忽略）。
\end{enumerate}
\onepicture[align=right, w=6.0cm]{figs/26.png}
%% answer:
\jdanswer{
\begin{enumerate}
	\item
	$T'=\left(\frac{r}{h}\right)^{3/2}T$

	\item
	$\tau=\frac{r^{3/2}}{\pi(h^{3/2}-r^{3/2})}\left(\arcsin\frac{R}{h}+\arcsin\frac{R}{r}\right)T$

\end{enumerate}
}
%% memo:
\memoanswer{
（1）设卫星 B 绕地心转动的周期为 $T'$，根据万有引力定律和圆周运动的规律有

$\frac{GMm}{h^2}=m\frac{4\pi^2}{T^2}h$

$\frac{GMm'}{r^2}=m'\frac{4\pi^2}{T'^2}r$

式中 $G$ 为引力常量，$M$ 为地球的质量，$m$、$m'$ 为卫星的质量

联立可得

$T'=\left(\frac{r}{h}\right)^{3/2}T$

（2）设卫星 A 和 B 连续不能直接通讯的时间间隔为 $\tau$

在此时间内，卫星 A 和 B 绕地心转动的角度分别为 $\alpha$ 和 $\alpha'$，则

$\alpha=\frac{\tau}{T}2\pi$

$\alpha'=\frac{\tau}{T'}2\pi$

若不考虑卫星 A 的公转，两卫星不能直接通讯时，卫星 B 的位置在图中 B 和 B' 点之间，图中内圆表示地球的赤道

由几何关系得 $\angle BOB'=2\left(\arcsin\frac{R}{h}+\arcsin\frac{R}{r}\right)$

当 $r<h$ 时，卫星 B 比卫星 A 转的快，考虑卫星 A 的公转后应有

$\alpha-\alpha'=\angle BOB'$

联立各式

$\tau=\frac{r^{3/2}}{\pi(h^{3/2}-r^{3/2})}\left(\arcsin\frac{R}{h}+\arcsin\frac{R}{r}\right)T$
}

\end{enumerate}

\end{document}
